\documentclass[12pt]{article}
\usepackage{amsmath,amssymb,amsthm}
\usepackage[margin=1in]{geometry}

\title{Project Euler Problem 739: Summation of Summations}
\author{}
\date{}

\begin{document}
\maketitle

\section{Problem Statement}
Take a sequence of length $n$. Discard the first term then make a sequence of partial sums. Repeat until one term remains, defining $f(n)$.

Starting with the Lucas-like sequence $1,3,4,7,11,18,29,47,\ldots$, find $f(10^8)\bmod 10^9+7$.

Given: $f(8)=2663$, $f(20)\equiv 742296999\pmod{10^9+7}$.

\section{Derivation}

\subsection{Linear Coefficients}
The final value is a linear combination:
\[
f(n) = \sum_{k=1}^{n} c_k \cdot a_k
\]
Since the first term is discarded at the first step, $c_1=0$. For $k\ge 2$:
\[
c_k = \binom{2(n-1)-k}{n-k}
\]

This can be shown by induction on the number of steps.

\subsection{Proof Sketch}
At each step, discarding position 1 and taking partial sums transforms coefficient vectors. After $n-1$ iterations on $n$ elements, the coefficient of the $k$-th original element involves a product of ``choosing paths'' through the triangular reduction, yielding the binomial coefficient formula.

\subsection{Computation}
For the Lucas sequence $L_k$ with $L_1=1$, $L_2=3$, $L_k=L_{k-1}+L_{k-2}$:
\[
f(n) = \sum_{k=2}^{n}\binom{2n-2-k}{n-k}\cdot L_k \pmod{10^9+7}
\]

\section{Algorithm}
\begin{enumerate}
\item Precompute factorials and inverse factorials modulo $10^9+7$ up to $2n$.
\item Iterate $k=2,\ldots,n$, computing $L_k$ and the binomial coefficient.
\item Accumulate the sum modulo $10^9+7$.
\end{enumerate}

\section{Complexity}
$O(n)$ time and space.

\section{Answer}

\[
\boxed{711399016}
\]

\end{document}
